% TOP_PROBLEM: {"id": 6700043, "title": "Scalar Curvature Question [?46]: But it is unclear if this remain true with \"area\" in place of \"length\"", "category_id": 6, "category": {"id": 6, "name": "geometry", "display_name": "Geometry", "description": "Euclidean and non-Euclidean geometry, geometric structures.", "slug": "geometry", "order_index": 6, "created_at": "2026-07-31T15:26:25.671Z"}, "status": "open"}
% FIRST_POSTED: null
% ENTRY_KIND: statement_only
% Imported from: batch15.tex; UnsolvedMath ID 6700043
\documentclass[11pt]{article}
\usepackage[a4paper,margin=25mm]{geometry}
\usepackage{amsmath,amssymb,mathtools}
\usepackage{fontspec,unicode-math}
\setmainfont{Latin Modern Roman}
\setsansfont{Latin Modern Sans}
\setmathfont{Latin Modern Math}
\usepackage{xurl,hyperref}
\hypersetup{colorlinks=true,linkcolor=blue,urlcolor=blue,pdftitle={UnsolvedMath Problem 6700043}}
\newcommand{\sref}[2]{\hyperlink{source:#1:#2}{[S#2]}}
\newcommand{\cataloguescope}[1]{\paragraph{Mathematical question.}#1}
\begin{document}
% UnsolvedMath ID 6700043; source catalogue code AMR-066-0043
\setcounter{section}{6700042}
\section{Scalar Curvature Question [?46]: But it is unclear if this remain true with "area" in place of "length"}\label{6700043}
{\small\sffamily UnsolvedMath ID 6700043 · Geometry}
\subsection{Definitions and mathematical statement}

\paragraph{Shared notation.}$\mathbb N=\{1,2,\ldots\}$, and $\mathbb Z,\mathbb Q,\mathbb R,\mathbb C$ denote the integers, rationals, reals and complex numbers. Any other conventions are specified locally.


Let $(Y,g_0)$ be an oriented Riemannian $n$-manifold. In the degree comparison consider oriented Riemannian $n$-manifolds $(X,g)$ and proper maps $f:X\to Y$ of nonzero degree. Proper means that inverse images of compact sets are compact; degree is the integer multiplying the orientation fundamental class (the relative class when the map sends boundary to boundary). A map is area non-increasing when the area of each smooth parametrized surface is not increased by its composition with $f$, with multiplicity counted. For a Lipschitz map, area is computed from its almost-everywhere differential. Scalar curvature $\operatorname{Sc}_g$ is the trace of the Ricci curvature tensor. The smooth area-extremality assertion here excludes all $C^2$ maps in this degree class satisfying both area non-increase and $\operatorname{Sc}_g(x)>\operatorname{Sc}_{g_0}(f(x))$ for every $x\in X$.
\paragraph{Statement recorded in the supplied source.}
If this smooth area-extremality assertion holds for $(Y,g_0)$, must it still hold when $C^2$ maps are replaced by Lipschitz maps in the same proper, nonzero-degree comparison class? \sref{6700043}{2}
\subsection{Short English statement}
If no smooth map of nonzero degree can decrease areas while sending an everywhere higher-scalar-curvature manifold onto the target, is the same exclusion valid for Lipschitz maps?
\subsection{Sources}
\begin{description}
\item[\hypertarget{source:6700043:1}{S1}] ULAM.AI and the UnsolvedMath Contributors, UnsolvedMath: A Curated Collection of Open Mathematics Problems, record 6700043 (AMR-066-0043). CC BY 4.0. \url{https://huggingface.co/datasets/ulamai/UnsolvedMath}.
\item[\hypertarget{source:6700043:2}{S2}] Misha Gromov, A Dozen Problems, Questions and Conjectures about Positive Scalar Curvature (18 October 2017), Section 4, printed pp. 10--14; arXiv:1710.05946v1. \url{https://arxiv.org/pdf/1710.05946v1}.
\end{description}
\label{end:6700043}
\end{document}
